\documentclass[12pt]{article}
% Préambule commun à tous les documents extraits de « La Revue CAPES »
\usepackage[T1]{fontenc}
\usepackage[utf8]{inputenc}
\usepackage{lmodern}
\usepackage{amsmath,amssymb,amsthm,amsfonts,mathrsfs}
\usepackage{setspace}
\usepackage{listings}
\usepackage{pgfplots}
\pgfplotsset{compat=1.18}
\usepackage{longtable}
\usepackage{adjustbox}
\usepackage{lettrine}
\usepackage{tkz-tab}
\usepackage{changepage}
\usepackage{pdflscape}
\usepackage{graphicx}
\usepackage{tikz}
\usepackage{xcolor}
\usepackage[a4paper,top=2cm,bottom=2.5cm,left=2.5cm,right=2cm]{geometry}
\usepackage{fancyhdr}
\usepackage{draftwatermark}
\SetWatermarkText{La RevueCapes}
\SetWatermarkLightness{0.9}
\SetWatermarkScale{2}
\usepackage{hyperref}
\hypersetup{colorlinks=true,linkcolor=blue!50!black,urlcolor=blue!50!black}

\definecolor{revue}{RGB}{20,60,120}

\pagestyle{fancy}
\fancyhf{}
\renewcommand{\headrulewidth}{0.4pt}
\setlength{\headheight}{15pt}

% \entete{Matière}{Titre}{Type de document}
\newcommand{\entete}[3]{%
  \fancyhead[L]{\small\textsc{La Revue CAPES} -- #1}%
  \fancyhead[R]{\small #2}%
  \fancyfoot[C]{\thepage}%
  \thispagestyle{plain}%
  \begin{center}
    {\color{revue}\rule{\textwidth}{1.2pt}}\\[4pt]
    {\small\textsc{Concours directs CAP-CEG / CAPES -- Burkina Faso}}\\[6pt]
    {\LARGE\bfseries\color{revue} #2}\\[6pt]
    {\large\itshape #1 -- #3}\\[2pt]
    {\color{revue}\rule{\textwidth}{1.2pt}}
  \end{center}
  \vspace{0.5cm}%
}

% Encadré pour les remarques ajoutées dans les corrigés
\newcommand{\remarque}[1]{\par\medskip\noindent\fcolorbox{revue}{revue!6}{\parbox{\dimexpr\linewidth-2\fboxsep-2\fboxrule}{\textbf{\color{revue}Remarque.} #1}}\par\medskip}


\begin{document}
\entete{Mathématiques}{Sujet type n°15}{Énoncé}
\begin{spacing}{1.5}

\medskip\textbf{Exercice 1}

1. Soit $q$ la forme quadratique définie par $q(x,y,z) = x(x-4y + 4z) +2yz$ et $A$ la matrice symétrique canoniquement associée à $q$.

(a) Déterminer la signature de q de deux manières :
 D'abord par la méthode de GAUSS, puis en étudiant le signe des mineurs principaux emboîtés de $A$.

(b) Donner la factorisation $QR$ de la matrice $A$.

2.(a) Étudier la convergence des séries suivantes : $\sum\frac{(n+2)^n}{n^{n+2}}$ et $\sum\frac{\sqrt{5n+1}}{n!}$

 (b) Calculer le rayon de convergence de la série entière suivante : $\sum \frac{3^n(n+7)^3x^n}{7^{n+3}}$.
 
3. Soit $f$ la fonction de $(\mathbb{R}_{+}^{*})^2$ dans $\mathbb{R}_{+}^{*}$, définie par $f(x,y)=\frac{2\sqrt{x}+3\sqrt{y}}{3x^2+xy+2y^2}$

 (a) Montrer que $f$ est positivement homogène, d'un degré $\alpha$ qu'on précisera.
 
 (b) Quel est le degré d'homogénéité de $(\frac{\partial^2f}{\partial x\partial y})^3$?
 
 (c) Calculer : $\frac{x^2}{f}\frac{\partial^2f}{\partial x^2}+\frac{2xy}{f}\frac{\partial^2f}{\partial x\partial y}+\frac{y^2}{f}\frac{\partial^2f}{\partial y^2}$.
 
\medskip\textbf{Exercice 2}

 Le service de dépannage d'un grand magasin dispose d'équipes intervenant sur appel de la clientèle. Pour des causes diverses, les interventions ont parfois lieu avec retard. On admet que les appels se produisent indépendamment les uns des autres, et que, pour chaque appel, la probabilité d'un retard est de 0,25.
 
 (1) Un client appelle le service à 4 reprises. On désigne par $X$ la variable aléatoire prenant pour valeurs le nombre de fois ou ce client a dû subir un retard.
 
 (a) Déterminer la loi de probabilité de $X$, son espérance, sa variance.
 
 (b) Calculer la probabilité de l'évènement : "Le client a au moins subi un retard".
 
 (2) Le nombre d'appels reçus par jour est une variable aléatoire $Y$ qui suit une loi de
 Poisson de paramètre $m$. On note $Z$ le nombre d'appels traités en retard.
 
 (a) Exprimer la probabilité conditionnelle de $Z = k$ sachant que $Y = n$.
 
 (b) En déduire la probabilité de $"Z = k$ et $Y = n"$.
 
 (c) Déterminer la loi de$ Z$. On trouvera que $Z$ suit une loi de Poisson de paramètre
 $m\times 0,25$.
 
 (3) En 2013, le standard a reçu une succession d'appels. On note $U$ le premier appel reçu en retard. Quelle est la loi de $U$? Quelle est son espérance?
 
 \medskip\textbf{Exercice 3 }
 
 Pour tout $x$ non nul, on pose
 
 $f(x)=\int_{x}^{2x}\frac{du}{u+\sin u}$
 
 On ne cherchera pas à calculer cette intégrale :
 
 1.a) Justifier l'existence de $f$
 
 b) Étudier la parité de $f$.
 
 c) Montrer que $\forall x>0, \hspace{0.1cm}f(x)>0$
 
 Dans la suite on supposera que $x>0$
 
 2. Calculer $f'(x)$, en déduire les variations de $f$
 
 3.(a) Trouver pour $u>1$, un encadrement de $\frac{1}{u+\sin u}$ puis pour $x>1$ un encadrement de $f(x)$.
 
(b) Déterminer \[
\lim_{x\to +\infty}f(x)
\]

\medskip
\textbf{Exercice 4}

Calculer la limite des suites suivantes :

1. \[
u_n=n\sum_{k=0}^{n-1}\frac{1}{n^2+k^2}
\]

2. \[
v_n=\prod_{k=1}^{n}\left(1+\frac{k^2}{n^2}\right)^{\frac{1}{n}}.
\]
 
 
 
 
 
 
 
 
 
 \newpage

\end{spacing}
\end{document}
